\chapter{Introduction}
\label{ch:intro}

%-------------------------------------------------------

% TODO: 1-2 paragraphs framing the assignment: what problem is being
% solved and why video brightness/contrast adjustment with histogram
% feedback is a standard image-processing task.

%%%%%%%%%%%%%%%%%%%%%%%%%%%%%%%%%%%%%%%%%%%%%%%%%%%%%
\section{Objectives}
\label{sec:objectives}
%%%%%%%%%%%%%%%%%%%%%%%%%%%%%%%%%%%%%%%%%%%%%%%%%%%%%

The assignment requires implementing a program that satisfies the
following eight tasks:

\begin{enumerate}
  \item Load and display the source video.
  \item Display the histogram of the source video.
  \item Display the same video as an adjustable preview.
  \item Display the histogram of the preview video.
  \item Loop both videos continuously, so that each restarts from its
        first frame as soon as it ends.
  \item Provide brightness and contrast scroll bars.
  \item Provide a save function, triggered by a button, a key, or a
        key combination.
  \item Persist the adjusted settings so that, when the program is
        reloaded, both windows reopen with the previous settings
        already applied.
\end{enumerate}


%%%%%%%%%%%%%%%%%%%%%%%%%%%%%%%%%%%%%%%%%%%%%%%%%%%%%
\section{Background}
\label{sec:background}
%%%%%%%%%%%%%%%%%%%%%%%%%%%%%%%%%%%%%%%%%%%%%%%%%%%%%

Image processing can be described as a function that takes one or
more images as input and produces an image as output. Once an image
is loaded, its pixel values can be manipulated in a variety of ways.
These manipulations fall into two broad categories: point operators
and neighborhood operators \citep{szeliski2022}. In a point operator,
the value of an output pixel depends only on the value of the
corresponding input pixel; in a neighborhood operator, the output
pixel also depends on the pixels surrounding it.

Two of the most common point operations are multiplication and
addition by a constant, combined into a single expression:
\begin{equation}
  g(x) = \alpha \cdot f(x) + \beta
  \label{eq:pointop-1d}
\end{equation}
where $f$ is the input image and $g$ is the output image. Applied to
every pixel of a two-dimensional image (see Equation
\ref{eq:pointop}), the same operation is written as
\begin{equation}
  g(x,y) = \alpha \cdot f(x,y) + \beta
  \label{eq:pointop}
\end{equation}

The parameter $\alpha$ is called the \emph{gain} and controls
contrast, while $\beta$ is called the \emph{bias} and controls
brightness. A dark image has many low-intensity pixels, so its
histogram shows a peak toward the left (low) end. Adding a constant
bias shifts every pixel by the same amount, which translates the
whole histogram to the right for $\beta > 0$ or to the left for
$\beta < 0$. Because pixel intensities are stored as 8-bit values,
any result above 255 is clipped to 255 and any result below 0 is
clipped to 0.


%%%%%%%%%%%%%%%%%%%%%%%%%%%%%%%%%%%%%%%%%%%%%%%%%%%%%
\section{Report Organization}
\label{sec:organization}
%%%%%%%%%%%%%%%%%%%%%%%%%%%%%%%%%%%%%%%%%%%%%%%%%%%%%

The remainder of this report is organized as follows. Chapter
\ref{ch:method} describes the development environment, the
application architecture, and how each of the eight tasks was
implemented. Chapter \ref{ch:results} presents sample outputs of the
program for each task. Chapter \ref{ch:conclusion} discusses the
results, notes the limitations encountered, and concludes the report.
